% !PRINTED 181
\weberpage{181}

\weq{9}{\varphi'(x_i) = 0 \qquad i = 1, 2, \ldots m,}

und wenn diese Determinante verschwindet, so giebt es (\S.\,24, III) ein Gr\"ossensystem

\weq{10}{x_1^{(0)}, x_2^{(0)} \ldots x_m^{(0)},}

dessen Elemente nicht alle verschwinden, und die, f\"ur die $x$ eingesetzt, den Gleichungen (9) gen\"ugen.

Nun hat man aber, wenn $x_i$, $x_i'$ irgend zwei Gr\"ossensysteme sind, und $\lambda$ ebenfalls eine willk\"urliche Gr\"osse bedeutet, die identischen Gleichungen

\weq{11}{\begin{aligned}
\varphi(x + \lambda x') &= \varphi(x) + \lambda \overset{i}{\Sigma} x_i \varphi'(x_i') + \lambda^2 \varphi(x')\\
\varphi(x') &= \tfrac{1}{2}\Sigma x_i' \varphi'(x_i'),
\end{aligned}}

[\S.\,16, (12) bis (15)], so dass, wenn f\"ur $x_i'$ die Werthe $x_i^{(0)}$ gesetzt werden

\weq{12}{\varphi(x_i + \lambda x_i^{(0)}) = \varphi(x),}

ganz unabh\"angig von $\lambda$ wird. Wenn wir nun annehmen, es sei $x_1^{(0)}$ von Null verschieden und wir setzen

\[
\lambda = -\frac{x_1}{x_1^{(0)}},
\]

ferner:

\weq{13}{\begin{aligned}
y_2 &= x_2 x_1^{(0)} - x_1 x_2^{(0)}\\
y_3 &= x_3 x_1^{(0)} - x_1 x_3^{(0)}\\
&\cdots\cdots\cdots\cdots\\
y_m &= x_m x_1^{(0)} - x_1 x_m^{(0)},
\end{aligned}}

so folgt

\weq{14}{x_1^{(0)\,2} \varphi(x) = \varphi(0, y_2, y_3 \ldots y_m),}

wodurch $\varphi(x)$ durch die $m - 1$ Ver\"anderlichen $y_2, y_3 \ldots y_m$ ausgedr\"uckt ist.

\begin{center}
\S.\,57.

\gesperrt{Transformation der quadratischen Formen in eine\\
Summe von Quadraten.}
\end{center}

Wir machen von dem zuletzt bewiesenen Satze sogleich eine wichtige Anwendung auf die Discriminante der Function (11) (als Function zweiten Grades von $\lambda$ betrachtet), wobei nicht
%CONT>>
